% ---------------------------------------------------------------------------
% Article template. Lines starting with % are comments: LaTeX ignores them.
% Change the text, compile (Ctrl+Enter), and the PDF on the right updates.
% ---------------------------------------------------------------------------
\documentclass[11pt]{article}

% Packages add features. amsmath: equations; graphicx: images.
\usepackage{amsmath}
\usepackage{graphicx}

% The title block. Separate several authors with \and.
\title{Your Paper Title}
\author{First Author \and Second Author}
\date{\today}

\begin{document}

\maketitle

% The abstract: one paragraph that says what you did and what you found.
\begin{abstract}
Summarise the problem, your approach and the main result in a few sentences.
\end{abstract}

% A section starts with \section{...}. Use \subsection for a level below.
% \label gives it a name you can refer to with \ref{sec:intro}.
\section{Introduction}\label{sec:intro}

Write your introduction here. To cite a source, add it to refs.bib and use its
key, for example \cite{knuth1984}. Section~\ref{sec:method} describes the method.

\section{Method}\label{sec:method}

Inline maths goes between dollar signs, like $E = mc^2$. A numbered equation
uses the equation environment:
\begin{equation}\label{eq:sum}
  S_n = \sum_{i=1}^{n} i = \frac{n(n+1)}{2}
\end{equation}
Refer to it as Equation~\eqref{eq:sum}.

% A figure. Upload an image to the project, then replace the box below with
% \includegraphics[width=0.8\linewidth]{your-image-file}
\begin{figure}[t]
  \centering
  \framebox[0.8\linewidth]{\rule{0pt}{3cm}Your figure here}
  \caption{Describe what the figure shows.}\label{fig:example}
\end{figure}

\section{Results}

Figure~\ref{fig:example} and Table~\ref{tab:example} show how to refer to
floats by their label.

% A table: one & between columns, \\ at the end of each row.
\begin{table}[t]
  \centering
  \begin{tabular}{lrr}
    \hline
    Method & Accuracy & Time (s) \\
    \hline
    Baseline & 0.81 & 12 \\
    Ours     & 0.87 & 9 \\
    \hline
  \end{tabular}
  \caption{A small results table.}\label{tab:example}
\end{table}

\section{Conclusion}

Close with what the results mean and what comes next.

% The bibliography is built from refs.bib.
\bibliographystyle{plain}
\bibliography{refs}

\end{document}
